%! Parametrically Defined Circles and Lines — Unit 8, Lesson 8.3 — Guided Notes
\documentclass[10pt]{article}
\usepackage{precalculus-article}
\usepackage{precalculus-boxes}
\newcommand{\TallMath}[1]{$\displaystyle #1\rule[-1.4em]{0pt}{3.2em}$}

\begin{document}

\pageheader{Unit 8, Lesson 8.3}{Guided Notes}
\namedateperiod

\vspace{0.08in}

%----------------------------------------------------------------------
% Objectives
%----------------------------------------------------------------------
\begin{objectivebox}
\textbf{I will be able to\ldots}
\begin{enumerate}[leftmargin=1.5em, itemsep=4pt]
  \item \textbf{LO 4.4.A} --- Express the motion of a point around a
    circle parametrically by applying transformations to
    $(x(t),\,y(t)) = (\cos t,\,\sin t)$, and parametrize the segment
    between two given points using an initial position and constant
    rates of change. \writeline
\end{enumerate}
\end{objectivebox}

\vspace{0.08in}

%----------------------------------------------------------------------
% Vocabulary
%----------------------------------------------------------------------
\begin{vocabbox}
\begin{description}[leftmargin=0pt, itemsep=6pt]
  \item[\termblanklong{unit circle parametrization}]
    The equations $x(t) = \cos t$, $y(t) = \sin t$, $0 \le t \le 2\pi$
    model one full \underline{\hspace{3cm}} revolution starting at
    \underline{\hspace{2cm}} and centered at the \underline{\hspace{2cm}}.

  \item[\termblanklong{radius / center}]
    In $(h + r\cos t,\; k + r\sin t)$: $r$ is the
    \underline{\hspace{2cm}} and $(h,\,k)$ is the
    \underline{\hspace{2cm}} of the circular path.

  \item[\termblanklong{orientation}]
    Using $(\cos t,\,\sin t)$ gives \underline{\hspace{3cm}} motion.
    Replacing $t$ with $-t$ gives \underline{\hspace{3cm}} motion.

  \item[\termblanklong{linear parametrization}]
    The parametrization $x(t) = x_1 + (x_2-x_1)t$,
    $y(t) = y_1 + (y_2-y_1)t$, $0\le t\le 1$, traces the segment from
    \underline{\hspace{2.5cm}} to \underline{\hspace{2.5cm}}.
\end{description}
\end{vocabbox}

\vspace{0.08in}

%----------------------------------------------------------------------
% Hook
%----------------------------------------------------------------------
\begin{hookbox}
\textbf{Entry Question:} A Ferris wheel has radius 30 ft with its center
35 ft above the ground. A rider boards at the rightmost seat.

\vspace{4pt}
\textbf{Q1.} What are the rider's starting coordinates $(x,\,y)$ relative
to a center at the origin?
\writeline

\textbf{Q2.} As the wheel turns counterclockwise, what do you expect
to happen to the rider's height $y$?
\writelines{2}
\end{hookbox}

\vspace{0.08in}

%----------------------------------------------------------------------
% Notes Box 1 — Unit Circle
%----------------------------------------------------------------------
\begin{notesbox}{LO 4.4.A --- Part 1: Unit Circle Parametrization (EK 4.4.A.1)}

\textbf{Definition:} The parametric function
$(x(t),\,y(t)) = (\cos t,\,\sin t)$, $0 \le t \le 2\pi$, traces the
\underline{\hspace{3cm}} circle exactly \underline{\hspace{2cm}} time,
moving \underline{\hspace{3cm}}.

\vspace{6pt}
\textbf{Complete the table:}

\vspace{4pt}
\begin{center}
\renewcommand{\arraystretch}{1.7}
\begin{tabular}{|c|c|c|c|}
\hline
$t$ & $x(t) = \cos t$ & $y(t) = \sin t$ & $(x,\;y)$ \\
\hline
$0$          & \blank{1.8cm} & \blank{1.8cm} & \blank{2.4cm} \\
\hline
$\pi/2$      & \blank{1.8cm} & \blank{1.8cm} & \blank{2.4cm} \\
\hline
$\pi$        & \blank{1.8cm} & \blank{1.8cm} & \blank{2.4cm} \\
\hline
$3\pi/2$     & \blank{1.8cm} & \blank{1.8cm} & \blank{2.4cm} \\
\hline
$2\pi$       & \blank{1.8cm} & \blank{1.8cm} & \blank{2.4cm} \\
\hline
\end{tabular}
\end{center}

\end{notesbox}

\vspace{0.08in}

%----------------------------------------------------------------------
% Notes Box 2 — Transformed Circles
%----------------------------------------------------------------------
\begin{notesbox}{LO 4.4.A --- Part 2: Transformations of the Circle (EK 4.4.A.2)}

\textbf{General form:}
\[
  x(t) = \underline{\hspace{1cm}} + \underline{\hspace{0.8cm}}\cos t, \qquad
  y(t) = \underline{\hspace{1cm}} + \underline{\hspace{0.8cm}}\sin t, \qquad
  0 \le t \le 2\pi
\]
\begin{itemize}[leftmargin=1.4em, itemsep=2pt]
  \item Center: \underline{\hspace{3cm}} \quad Radius: \underline{\hspace{2cm}}
  \item Counterclockwise when $t$ increases \quad / \quad Clockwise: replace $t$
    with \underline{\hspace{1.5cm}}
\end{itemize}

\vspace{8pt}
\textbf{Ferris Wheel Example.} Radius 30 ft, center at $(0,\,35)$, counterclockwise:

\vspace{4pt}
$x(t) = $ \blank{4cm} \qquad $y(t) = $ \blank{4cm}

\vspace{6pt}
Evaluate at $t = \pi/2$: \quad $x(\pi/2) = $ \blank{2cm}, \quad $y(\pi/2) = $ \blank{2cm}

The rider is at position \blank{3cm} after a quarter turn.

\vspace{6pt}
\textbf{Clockwise version} (replace $t$ with $-t$):

$x(t) = $ \blank{4cm} \qquad $y(t) = $ \blank{4cm}

\end{notesbox}

\vspace{0.08in}

%----------------------------------------------------------------------
% Notes Box 3 — Linear Segment
%----------------------------------------------------------------------
\begin{notesbox}{LO 4.4.A --- Part 3: Linear Parametrization (EK 4.4.A.3)}

\textbf{Segment from $(x_1,\,y_1)$ to $(x_2,\,y_2)$:}
\[
  x(t) = x_1 + (x_2-x_1)\,t, \qquad y(t) = y_1 + (y_2-y_1)\,t, \qquad 0 \le t \le 1
\]
At $t=0$: position is \underline{\hspace{2.5cm}}.
At $t=1$: position is \underline{\hspace{2.5cm}}.
At $t=0.5$: position is the \underline{\hspace{2cm}}.

\vspace{6pt}
\textbf{Example:} Segment from $(2,\,5)$ to $(6,\,1)$.

\vspace{4pt}
$x(t) = 2 + \bigl($ \blank{1.2cm} $\bigr)t = $ \blank{3.2cm} \qquad
$y(t) = 5 + \bigl($ \blank{1.2cm} $\bigr)t = $ \blank{3.2cm}

\vspace{6pt}
\textbf{Complete the table:}

\vspace{4pt}
\begin{center}
\renewcommand{\arraystretch}{1.6}
\begin{tabular}{|c|c|c|c|}
\hline
$t$ & $x(t)$ & $y(t)$ & $(x,\;y)$ \\
\hline
$0$   & \blank{1.8cm} & \blank{1.8cm} & \blank{2.4cm} \\
\hline
$0.5$ & \blank{1.8cm} & \blank{1.8cm} & \blank{2.4cm} \\
\hline
$1$   & \blank{1.8cm} & \blank{1.8cm} & \blank{2.4cm} \\
\hline
\end{tabular}
\end{center}

\end{notesbox}

\end{document}
